\section{Criterio de selección}
\label{sec:candidate-selection}

\HAFOCardCandidates

Seleccionamos sistemas en los que la representación algebraica y el
problema dinámico puedan contrastarse por separado. Consideramos:

\begin{enumerate}
 \item derivada ordinaria o Caputo conmensurable \(q\in(0,1]\);
 \item ecuaciones, parámetros, orden e iniciales publicados;
 \item no linealidades representables por D0, T1, A1 o M1;
 \item un resultado bibliográfico sobre atractores ocultos, cuencas o coexistencia;
 \item tamaño compatible con ABM de memoria completa y una reproducción
 independiente;
 \item una estructura que permita evaluar realizabilidad, memoria o coexistencia.
\end{enumerate}

La búsqueda principal excluye sistemas inconmensurables, Caputo--Fabrizio,
Atangana--Baleanu, retardados, discretos o complejos, cuyos símbolos de
transferencia, condiciones de estabilidad y protocolos de memoria difieren de
Caputo conmensurable.

\section[Sistema Muñoz--Pacheco sin equilibrio]
{Sistema fraccionario sin equilibrio de Muñoz--Pacheco}
\label{sec:cand-munoz}

La fuente primaria \cite{munoz2018} propone
\begin{equation}
\begin{aligned}
 {}^{\mathrm C}D^q x&=yz+x(y-a),\\
 {}^{\mathrm C}D^q y&=1-|x|,\\
 {}^{\mathrm C}D^q z&=-xy-z.
\end{aligned}
\label{eq:cand-munoz}
\end{equation}
Para \(a=0.35,q=0.97\), publica un atractor caótico desde
\((1,1,1)\). Para \(a=0.35,q=0.996\), publica coexistencia de un candidato
caótico desde \((1,1,1)\) y uno periódico desde \((0,75,-50)\), junto con un
corte de cuenca. Estos datos permiten comparar localización, coexistencia y
dependencia del orden fraccionario.

\subsection{Realización M1 exacta}

Con \(X=(x,y,z)^{\mathsf T}\),
\[
 \Phi(X)=(yz,xy,|x|)^{\mathsf T},
\]
se tiene
\begin{equation}
 {}^{\mathrm C}D^qX=A X+d+B\Phi(X),
 \label{eq:cand-munoz-mimo}
\end{equation}
donde
\begin{equation}
 A=\begin{pmatrix}
 -a&0&0\\0&0&0\\0&0&-1
 \end{pmatrix},\quad
 d=\begin{pmatrix}0\\1\\0\end{pmatrix},\quad
 B=\begin{pmatrix}
 1&1&0\\
 0&0&-1\\
 0&-1&0
 \end{pmatrix}.
 \label{eq:cand-munoz-matrices}
\end{equation}
La multiplicación da
\[
 AX+d+B\Phi=
 \begin{pmatrix}
 -ax+yz+xy\\1-|x|\\-z-xy
 \end{pmatrix},
\]
idéntica a \cref{eq:cand-munoz}.

Con \(C=I\) y \(z_f=s^q\), el bloque lineal es diagonal:
\[
 z_fI-A=\operatorname{diag}(z_f+a,z_f,z_f+1).
\]
Para \(z_f\notin\{-a,0,-1\}\), invertimos cada entrada diagonal y
multiplicamos por las tres columnas de \(B\). Obtenemos
\begin{equation}
 W(z_f)=(z_fI-A)^{-1}B
 =\begin{pmatrix}
 \dfrac1{z_f+a}&\dfrac1{z_f+a}&0\\[0.8em]
 0&0&-\dfrac1{z_f}\\[0.8em]
 0&-\dfrac1{z_f+1}&0
 \end{pmatrix}.
 \label{eq:cand-munoz-transfer}
\end{equation}
La transferencia describe la respuesta a los canales no lineales con dato
inicial nulo; el término constante \(d\) interviene por separado en el
balance DC. Si los estados contienen un primer armónico, los productos
\(xy\) y \(yz\) producen una componente media y un segundo armónico.
Por ello tomamos M1 con \(H\ge2\) como primera aproximación. El valor
absoluto puede generar armónicos de órdenes superiores: se integra por
regiones en el cálculo de Fourier, o se aproxima mediante una homotopía
suavizada cuyo error debe controlarse.

\subsection{Demostración completa de ausencia de equilibrios}

En equilibrio, la segunda ecuación exige \(|x|=1\), de modo que
\(x=\pm1\neq0\). La tercera da
\[
 z=-xy.
\]
Al sustituir en la primera,
\[
 0=y(-xy)+x(y-a)
 =x(-y^2+y-a).
\]
Como \(x\neq0\), debe cumplirse
\begin{equation}
 y^2-y+a=0.
 \label{eq:cand-munoz-equilibrium}
\end{equation}
Su discriminante es
\[
 \Delta_y=1-4a.
\]
Para \(a=0.35\), \(\Delta_y=-0.4<0\); por tanto, no existe equilibrio real.
Bajo la definición usual,
cualquier atractor del sistema es oculto por ausencia de equilibrios. Aun así,
deben reproducirse recurrencia y caos: la ausencia de equilibrios no transforma
una trayectoria transitoria o divergente en atractor.

Para \(x\neq0\), el Jacobiano regional es
\begin{equation}
 J(X)=
 \begin{pmatrix}
 y-a&z+x&y\\
 -\operatorname{sign}(x)&0&0\\
 -y&-x&-1
 \end{pmatrix}.
 \label{eq:cand-munoz-jacobian}
\end{equation}
La ausencia de equilibrios hace inaplicable el criterio de Matignon local. La
validación se concentra en el integrador no suave, la coexistencia y la
robustez de los diagnósticos.

\begin{proposalbox}
Para comparar las respuestas publicadas mantenemos \(a=0.35\), los órdenes
\(q=0.97\) y \(q=0.996\), las condiciones iniciales indicadas y la misma
definición de \(|x|\). La sucesión de truncaciones \(H=2,3,4\) permite
examinar por separado el balance DC, la fase, los residuos retenidos y la
cola física \(\epsilon_H\) de \cref{eq:th-physical-harmonic-tail};
\(\eta_H\) mide la contribución del diccionario elegido. La condición inicial
reconstruida define después un PVI de Caputo con memoria completa. El corte
\(x=0\), las respuestas coexistentes y la frontera de escape son los objetos
dinámicos de comparación. Al no existir equilibrios, el análisis de cuenca
reemplaza al sondeo de sus vecindades; la concordancia entre ABM y un segundo
integrador se evalúa sobre el mismo PVI.
\end{proposalbox}

\section{Jerk cuadrático entero con equilibrio estable}
\label{sec:cand-jerk}

Liu, Sang, Wang y Ahmad \cite{liu2021} estudian
\begin{equation}
 \dot x=y,\qquad
 \dot y=z,\qquad
 \dot z=-x-y-1.1z-2.59z^2+xy.
 \label{eq:cand-jerk}
\end{equation}
La fuente reporta un atractor caótico oculto desde
\((0,-1.5,0)\), además de cortes de cuenca y una realización de circuito.

\subsection{Parte lineal, equilibrio y estabilidad}

Con \(X=(x,y,z)^{\mathsf T}\),
\begin{equation}
 A=\begin{pmatrix}
 0&1&0\\0&0&1\\-1&-1&-1.1
 \end{pmatrix}.
 \label{eq:cand-jerk-A}
\end{equation}
Las dos primeras ecuaciones de equilibrio dan \(y=z=0\); la tercera da
\(x=0\). El origen es el único equilibrio. Su polinomio característico es
\begin{equation}
 p(s)=s^3+1.1s^2+s+1.
 \label{eq:cand-jerk-charpoly}
\end{equation}
En efecto, al expandir el determinante de
\[
 sI-A=\begin{pmatrix}s&-1&0\\0&s&-1\\1&1&s+1.1\end{pmatrix}
\]
por la primera fila resulta \(s[s(s+1.1)+1]+1\).
La primera columna de la tabla de Routh es
\[
 1,\qquad 1.1,\qquad \frac{1.1\cdot1-1}{1.1},\qquad 1.
\]
Sus entradas son positivas. El equilibrio único es
asintóticamente estable. Un atractor coexistente no puede localizarse
siguiendo la variedad inestable del origen y, si su cuenca está separada de una
vecindad del mismo, es oculto.

\subsection{Realización M1 y obstrucción escalar}

Definimos
\[
 \Phi(X)=(xy,z^2)^{\mathsf T},\qquad
 B=\begin{pmatrix}0&0\\0&0\\1&-2.59\end{pmatrix}.
\]
Entonces
\[
 AX+B\Phi(X)=
 (y,z,-x-y-1.1z+xy-2.59z^2)^{\mathsf T},
\]
que reproduce exactamente el campo.

Aunque los dos canales actúan en la dirección \(e_3\), no pueden reducirse a
una no linealidad cuadrática de un único observable. Para
\[
 Q(x,y,z)=xy-2.59z^2
\]
se tiene
\begin{equation}
 \nabla^2Q=
 \begin{pmatrix}
 0&1&0\\1&0&0\\0&0&-5.18
 \end{pmatrix},
 \qquad \operatorname{rank}\nabla^2Q=3.
 \label{eq:cand-jerk-hessian}
\end{equation}
El determinante de esta Hessiana es \(5.18\ne0\), de donde procede el
rango tres. En cambio, \(Q=\kappa(r^{\mathsf T}X)^2\) implicaría
\(\nabla^2Q=2\kappa rr^{\mathsf T}\), de rango a lo sumo uno. Esto falsifica
D0 y justifica M1.

Para obtener la transferencia resolvemos \((sI-A)v=e_3\). Las dos
primeras filas dan \(v_2=sv_1\), \(v_3=s^2v_1\), y la última da
\(p(s)v_1=1\). Como la segunda columna de \(B\) es \(-2.59e_3\),
para \(p(s)\ne0\) obtenemos
\begin{equation}
 W(s)=\frac1{p(s)}
 \begin{pmatrix}
 1&-2.59\\
 s&-2.59s\\
 s^2&-2.59s^2
 \end{pmatrix}.
 \label{eq:cand-jerk-transfer}
\end{equation}
Los términos cuadráticos requieren DC y segundo armónico. Este sistema ofrece
un control entero para evaluar si M1 amplía la cobertura respecto de los casos
escalares.

\section{Sistema Liu extendido fraccionario}
\label{sec:cand-extended-liu}

La fuente \cite{wang2020} usa
\begin{equation}
\begin{aligned}
 {}^{\mathrm C}D^q x&=a(y-x),\\
 {}^{\mathrm C}D^q y&=bx-kxz+d\,\operatorname{sign}(w)+m,\\
 {}^{\mathrm C}D^q z&=-cz+hx^2,\\
 {}^{\mathrm C}D^q w&=-gy.
\end{aligned}
\label{eq:cand-liu}
\end{equation}
Una realización exacta usa
\[
 \Phi=(xz,x^2,\operatorname{sign}w)^{\mathsf T},
\]
\[
 A=\begin{pmatrix}
 -a&a&0&0\\ b&0&0&0\\0&0&-c&0\\0&-g&0&0
 \end{pmatrix},\quad
 d_0=(0,m,0,0)^{\mathsf T},
\]
\[
 B=\begin{pmatrix}
 0&0&0\\-k&0&d\\0&h&0\\0&0&0
 \end{pmatrix}.
\]
La multiplicación permite comprobar cada canal:
\[
 AX+d_0+B\Phi=
 \begin{pmatrix}
 -ax+ay\\bx-kxz+d\operatorname{sign}(w)+m\\-cz+hx^2\\-gy
 \end{pmatrix}.
\]
Así, el término constante y el salto del signo permanecen explícitos y no
se absorben en una linealización.

Para \(a=10,b=40,c=2.5,k=1,h=4,g=5,d=m=20\), las ecuaciones de equilibrio
dan sucesivamente \(y=x\), \(y=0\), \(x=0\), \(z=0\), y
\[
 20\,\operatorname{sign}(w)+20=0.
\]
Con \(\operatorname{sign}(0)=0\), la solución es la semirrecta
\begin{equation}
 x=y=z=0,\qquad w<0,
 \label{eq:cand-liu-equilibria}
\end{equation}
no un sistema sin equilibrios. Su evaluación requiere fijar \(q\), la rama
dinámica, las condiciones iniciales, la convención en \(w=0\) y el protocolo
para una variedad continua de equilibrios.

\HAFOStudyNeedspace
\section{Candidatos adicionales y controles}
\label{sec:candidate-queue}

\begin{longtable}{@{}L{0.23\linewidth}L{0.20\linewidth}L{0.24\linewidth}L{0.25\linewidth}@{}}
\caption{Familias compatibles y problemas dinámicos asociados.}
\label{tab:candidate-portfolio}\\
\toprule
Familia & Ruta estructural & Construcción disponible & Problema por resolver\\
\midrule
\endfirsthead
\toprule
Familia & Ruta estructural & Construcción disponible & Problema por resolver\\
\midrule
\endhead
Leonov--Kuznetsov (162) & D0+A1 & Ecuaciones y semilla parciales &
Precisar la característica PWA y el criterio por etapa.\\
Genesio--Tesi & D0 sesgado, \(H\ge2\) & Realización
\(\psi(x)=x^2\) & Fijar parámetros e iniciales para el estudio dinámico.\\
Jerk con seno hiperbólico & D0, DF con \(I_1\) & Derivación simbólica &
Resolver ambigüedad de fuente \(\sinh\) frente a función inversa.\\
Oscilador memristivo & T1 sobre variedad & Cierre cúbico sesgado &
Constante de la variedad, parámetros e iniciales.\\
Jerk PWA & D0 PWA & Estructura tipo Chua & Fijar el campo regional y evaluar
el efecto de sus superficies de conmutación.\\
Fischer jerk exponencial & D0 con DC & Parámetros y diagnóstico &
Control bibliográfico/Lyapunov; clasificación de cuenca no reproducida
localmente.\\
Lorenz generalizado & M1, \(H\ge2\) & Álgebra completa verificada &
Raíz armónica, continuación, integración y cuenca.\\
Rabinovich--Fabrikant & M1, \(H\ge3\) & Álgebra completa verificada &
Raíz armónica, continuación, integración y cuenca.\\
Sistema financiero Fischer & M1+PWA & Banco Lyapunov & Resolver discrepancias
y distinguir la no linealidad \(|x|\) de otros modelos financieros.\\
Four-wing Fischer & M1 cuadrático & Diagnósticos de Lyapunov & Localizar una
semilla y estudiar la cuenca.\\
Debbouche fraccionario & M1+signo & DOI y parámetros parciales
\cite{debbouche2021} & Fijar ecuaciones, órdenes e iniciales de una rama.\\
\bottomrule
\end{longtable}

\section{Relación entre los problemas de localización}
\label{sec:candidate-order}

El jerk cuadrático entero aísla la extensión D0 \(\to\) M1 sin introducir
el costo de la memoria. Muñoz--Pacheco añade tres dificultades distintas:
la no suavidad, la historia de Caputo y la coexistencia dependiente del orden.
Lorenz generalizado y RF permiten estudiar el balance multiharmónico con
realizaciones polinómicas explícitas, mientras que Genesio--Tesi proporciona
un contraste escalar con sesgo. Estas diferencias hacen posible atribuir un
fallo de localización a una etapa concreta, en lugar de agruparlo bajo la
etiqueta de complejidad del sistema.

Liu extendido requiere además definir la solución en la discontinuidad y
analizar una variedad continua de equilibrios. En los modelos de Debbouche,
la elección de una rama debe fijar simultáneamente ecuaciones, órdenes y
condiciones iniciales. Las realizaciones algebraicas anteriores permiten
formular esas comparaciones, pero el destino de cada semilla se determina
integrando el correspondiente problema dinámico.
